\section{Experimental Methodology}
\subsection{Power Measurements}
A node was powered at 4.2~V (fully charged cell) from a Rohde \& Schwarz HMP2030 laboratory supply, which measures current down to 100~$\mu$A. Firmware cycled the node through four states in a loop, holding each for 10~s so that the reading could stabilize: deep sleep (everything off, only the RTC interrupt enabled), awake with sensors on, LoRa transmit/receive, and half-sleep (sensors off, MCU in light sleep, LoRa radio listening). Three complete loops were observed to check repeatability.

\subsection{Field Tests}
Testing proceeded in stages, in 2025:
\begin{itemize}
\item \textit{Benchtop test}, 1--4 June: four nodes (HW\_1, HW\_3, HW\_5, HW\_7) uploaded roughly hourly to ThingSpeak, giving 72 records per node, and mesh relaying was observed through debug output.
\item \textit{Preliminary Test I}, 8--29 June, and \textit{Preliminary Test II}, 1 August--6 October: four single-cell nodes on turf roofs on the University of Iceland campus, within about 50--100~m of a base station in an office. Test II used larger, better antennas. Nearby Icelandic Meteorological Office stations 1469 and 1477 served as context for the measurements.
\item \textit{LoRa range test} in Vellir, Hafnarfj{\"o}r{\dh}ur: the base station was mounted on a fence at chest height and sent randomly encoded packets once per second while a node in a vehicle was driven around the area. RSSI was recorded continuously and received packets were checked for corruption. Only bandwidth and spreading factor were varied; the other parameters were fixed.
\item \textit{Keldur deployment} from 17 October: four nodes with three cells each were buried 60--85~mm below the turf surface, shallower than the original approximately 110~mm design target, in a soil layer about 170~mm deep, with one node installed beneath each turf roof. The base station sat in the attic of one turf house next to a Wi-Fi router. Nodes 1 and 3 had the air sensor inside the enclosure, and nodes 5 and 7 had it mounted externally above the enclosure, to compare the two arrangements. The nodes were not visited or serviced afterwards. Measurements were compared with the Sámsstaðir automatic station (6222), 9.75~km away.
\end{itemize}

\subsection{Soil Moisture Calibration}
An Icelandic soil sample was oven-dried at 105~$^\circ$C for 32 hours and cooled. Water was then added in 10\% increments relative to the dry mass, up to 40\%, with multiple readings after each increment. The capacitive sensors were installed side by side with a HOBO MX2307 logger, whose sensor is documented as a METER TEROS 10, as reference instrument \cite{onset_mx2307}, and the gravimetric moisture content was taken as the actual value (Eq.~\ref{eq:gravimetric}). The procedure was run twice, with a 2069.5~g sample on 8 November and a 400~g sample on 16 December. A second-order polynomial relating ADC value to gravimetric moisture was fitted to the mean of the two sensors in Test 1, and its root-mean-square error (RMSE) against the gravimetric values was computed.
\begin{equation}
\text{MC}\,(\%) = \frac{m_{wet}-m_{dry}}{m_{dry}}\times 100
\label{eq:gravimetric}
\end{equation}
